\chapter{Capitolo Esempio}
\label{chap:Capitolo2}

Lorem ipsum dolor sit amet, consectetur adipisci elit, sed do eiusmod tempor incidunt ut labore et dolore magna aliqua. Ut enim ad minim veniam, quis nostrum exercitationem ullamco laboriosam, nisi ut aliquid ex ea commodi consequatur. Duis aute irure reprehenderit in voluptate velit esse cillum dolore eu fugiat nulla pariatur. Excepteur sint obcaecat cupiditat non proident, sunt in culpa qui officia deserunt mollit anim id est laborum.

In questo capitolo andremo a discutere ...
  
\section{Sezione Esempio}
\label{sec:real-time}
Quello in Figura \ref{fig:rocker}  (esempio di riferimento a figura) ...

% \begin{figure}[htpb!]
%   \centering
%   \includegraphics[width=0.5\textwidth]{Rockerduck}
%   \caption{Esempio di figura}
%   \label{fig:rocker}
% \end{figure}

Esempio elenco puntato ...
\begin{itemize}
	\item item 1
	\item item 2
	\item item 3
\end{itemize}


\section{Section2}

\subsection{Subsection Esempio}
\label{sec:handshake}

\subsection{Subsection Esempio}
\label{sec:handshake}

\begin{lstlisting}[caption={Esempio di listing}, style=javaScriptCode]
	GET /chat HTTP/1.1
	Host: server.example.com
	Upgrade: websocket
	Connection: Upgrade
	Sec-WebSocket-Key: dGhlIHNhbXBsZSBub25jZQ==
	Origin: http://example.com
	Sec-WebSocket-Protocol: chat, superchat
	Sec-WebSocket-Version: 13
\end{lstlisting} 


\begin{table}[htbp]
	\begin{center}
		\begin{tabular}{|l|l|l|l|l|l|}
			\hline
			Versione & Chrome & Firefox & Internet Explorer & Opera & Safari \\
			\hline
			76 & 6 & 4.0 & No & 11.00(disabilitato) & 5.0.1\\
			\hline
			7 & No & 6.0 & No & No & No \\
			\hline
			10 & 14 & 7.0 & HTML5 Labs & ? & ?\\
			\hline
			RFC 6455 & 16 & 11.0 & 10 & 12.10 & 6.0\\
			\hline
		\end{tabular}
	\end{center}
	\caption{Esempio di Tabella}
	\label{tab:browser}
\end{table}

\begin{table}[htbp]
	\begin{center}
		\begin{tabular}{|l|l|l|l|l|l|}
			\hline
			Versione & Android & Firefox Mob. & IE Mob. & Opera Mob. & Safari Mob.\\
			\hline
			76 & ? & ? & ? & ? & ?\\
			\hline
			7 & ? & ? & ? & ? & ? \\
			\hline
			10 & ? & 7.0 & ? & ? & ?\\
			\hline
			RFC 6455 & 16(Chrome) & 11.0 & ? & 12.10 & 6.0\\
			\hline
		\end{tabular}
	\end{center}
	\caption{Esempio di Tabella}
	\label{tab:mobile}
\end{table}
Nelle Tabelle \ref{tab:browser} e \ref{tab:mobile} è possibile vedere, rispettivamente per desktop e per mobile, il supporto dei vari browser per le diverse specifiche delle WebSocket.
Il codice completo dell'esempio è disponibile sul mio GitHub\footnote{\url{https://github.com/Glydric/TesiTriennale}} (Esempio di link).

\section{Qui}


Ciao Ciao

\newpage
Ciao Ciao


Ciao Ciao
